\chapter{流形学习}
\label{chap:manifold}
\begin{quote}\small
\textit{本章摘要。}本章研究高维观测中存在低维几何结构时，如何构造可解释的低维表示。内容从流形假设、内在维度和局部邻域开始，依次说明 Isomap、局部线性嵌入、拉普拉斯特征映射和扩散映射的机制，再讨论 t-SNE 与 UMAP 的可视化目标。最后给出方法选择、采样条件和结果解释的检查原则。低维图形描述的是所选几何目标，不能单独证明类别可分或预测有效。
\end{quote}

上一章讨论联邦学习如何在数据分散时训练模型。本章转向另一个问题。设备传感器可能记录几十个变量，但运行状态主要由负载、转速和磨损程度等少数连续因素决定。样本因此可能集中在高维空间的一条曲线或一个曲面附近。PCA 能找到一个最好的线性子空间，但折叠、弯曲或周期结构会使线性投影产生较大的失真。流形学习利用样本邻域推断局部几何，再按照不同目标构造低维坐标。

本章只讨论无监督表示和可视化。标签若参与目标函数，任务就变成监督降维或预测模型，评价标准也会改变。正文沿着“局部结构如何变成全局距离或谱坐标”的线索组织方法。Isomap保留图上的全局测地距离。LLE保留每个点由近邻线性重构的权重。拉普拉斯特征映射和扩散映射使相邻样本在低维空间中平滑。t-SNE与UMAP主要用于显示邻域关系，远距离和坐标轴不具有统一的物理解释。

\section{流形假设与内在维度}

\subsection{由高维观测产生的低维结构}
设第$i$个样本的观测为$\bm x_i\in\mathbb R^d$。流形假设认为，存在一个远小于$d$的内在维度$r$，以及局部光滑的映射$g$，使
\begin{equation}
 \bm x_i=g(\bm z_i)+\bm\varepsilon_i,
 \qquad \bm z_i\in\mathbb R^r,\quad 1\le r\ll d .
 \label{eq:manifold-generative}
\end{equation}
$\bm z_i$是潜在坐标，$g:\mathbb R^r\to\mathbb R^d$把状态因素映射到传感器空间，$\bm\varepsilon_i$是测量噪声或未建模变化。理想无噪声集合$\mathcal M=g(\mathcal Z)$称为$r$维流形的嵌入。局部光滑且雅可比矩阵$Dg(\bm z)$满列秩时，附近样本可以用$r$个局部坐标描述。满秩条件排除了局部塌缩。局部单射条件避免两个不同状态在一个小邻域内被混合。

内在维度描述产生数据所需的局部自由度。它不等于观测变量的个数，也不等于某种算法输出的坐标数。圆周样本$(\cos t,\sin t)$只有一个内在维度，却需要两个观测坐标表示整条闭合曲线。若测量噪声与流形法向方向相当，算法看到的将是有厚度的管状点云，局部线性近似会变得不稳定。流形假设还要求采样覆盖足够均匀。只有曲线的一小段被观测到时，算法无法判断它属于直线、圆周还是更高维曲线。

图\ref{fig:manifold-folded-neighborhood}区分了局部欧氏距离和沿流形的距离。折叠曲线两段在观测空间中很近，但沿曲线行走需要经过较长路径。邻域半径过大时会把两段错误连接，邻域过小时图可能断开。
\begin{figure}[htbp]
\centering
\begin{tikzpicture}[x=0.9cm,y=0.72cm]
 \draw[gray!30] (0,0) grid (7,4.2);
 \draw[blue!75!black,very thick] plot[smooth] coordinates {(0.4,0.7)(1.2,1.5)(2.1,2.8)(3.2,3.4)(4.2,2.7)(5.0,1.3)(6.1,0.8)};
 \foreach \x/\y in {0.4/0.7,1.2/1.5,2.1/2.8,3.2/3.4,4.2/2.7,5.0/1.3,6.1/0.8}{\fill[red!70!black] (\x,\y) circle (2pt);}
 \draw[orange!80!black,dashed] (2.1,2.8) circle (0.48cm);
 \draw[orange!80!black,dashed] (4.2,2.7) circle (0.48cm);
 \draw[<->,thick] (2.1,2.35)--(4.2,2.35) node[midway,below,font=\scriptsize]{欧氏距离较近};
 \draw[->,thick,gray!75!black] (2.2,3.05) .. controls (3.1,4.1) and (4.0,4.1) .. (4.1,3.0) node[midway,above,font=\scriptsize]{沿流形的路径};
\end{tikzpicture}
\caption{折叠曲线上的两类距离。局部邻域应沿曲线取点，过大的邻域会在折叠处形成捷径。图为定性示意。}
\label{fig:manifold-folded-neighborhood}
\end{figure}

\subsection{度量、采样与噪声}
欧氏距离定义为$d_E(\bm x_i,\bm x_j)=\|\bm x_i-\bm x_j\|_2$。它适合各坐标具有可比尺度且局部曲率较小的情形。若传感器单位不同，应只用训练集均值和尺度进行标准化。对零方差特征需要删除或预先规定处理方式。局部邻域可用$k$近邻或半径$\rho$构造。$k$过小会导致图断连和噪声边主导，$k$过大则可能跨越折叠。应检查连通分量、近邻边长度和对$k$的稳定性。

流形方法依赖有限样本近似连续结构。样本应覆盖关心的区域，局部采样间距应小于所要辨识的曲率尺度。密度急剧变化时，固定$k$会在稀疏区使用很大的物理半径，在密集区使用很小的半径。噪声会制造错误近邻，也会使局部权重矩阵病态。常见处理包括增加邻域、使用距离核、去噪或对局部矩阵加小的对角线正则项。任何处理都会改变被保持的几何对象，不能把它当作无影响的预处理。

\section{PCA的线性限制与图表示}
PCA在中心化数据矩阵$X\in\mathbb R^{n\times d}$上寻找正交矩阵$V\in\mathbb R^{d\times q}$，通过
\begin{equation}
 \min_{V^{\mathsf T}V=I_q}\|X-XVV^{\mathsf T}\|_F^2
 \label{eq:manifold-pca-baseline}
\end{equation}
最小化线性投影的平均平方重构误差。上一章已经说明它等价于保留协方差矩阵的前$q$个特征方向。该目标适合点云接近平面或线性子空间的情形。对弯曲曲线，投影会把相距很远的状态压到相近坐标。对圆周，单个线性坐标无法保持闭合结构。核PCA可在特征空间中弯曲表示，但核尺度决定邻域和全局几何，且原空间重建需要额外的原像假设\cite{scholkopf1998kpca}。

为了让非线性方法访问局部结构，定义无向邻域图$G=(\mathcal V,\mathcal E)$。节点集合$\mathcal V=\{1,\ldots,n\}$，若样本$i,j$互为近邻或满足半径条件，则$(i,j)\in\mathcal E$。图\ref{fig:manifold-neighborhood-graph}中的实线只连接局部近邻，Isomap沿这些边累加路径长度，LLE用这些邻居求局部重构权重，图谱方法则用边权惩罚相邻节点的坐标差异。图中的虚线捷径没有局部几何依据。若把它加入图中，Isomap的最短路径会被缩短，谱方法和扩散方法也会把折叠两侧错误地视为相近区域。边长记为$\ell_{ij}=\|\bm x_i-\bm x_j\|_2$。若算法需要相似度，可定义
\begin{equation}
 w_{ij}=\exp\left(-\frac{\ell_{ij}^2}{4t}\right),\qquad t>0,
 \label{eq:manifold-heat-weight}
\end{equation}
并令未连边权重为$0$。$\ell_{ij}$是距离，$w_{ij}$是相似度。两者在后续目标中承担不同作用。

\begin{figure}[htbp]
\centering
\begin{tikzpicture}[x=1.0cm,y=0.8cm]
 \foreach \p/\name in {(0.5,1)/1,(1.4,1.55)/2,(2.4,1.0)/3,(3.4,1.7)/4,(4.4,1.1)/5}{\fill[blue!70!black] \p circle (2.5pt) node[above=3pt,font=\scriptsize]{$\name$};}
 \draw[thick] (0.5,1)--(1.4,1.55)--(2.4,1.0)--(3.4,1.7)--(4.4,1.1);
 \draw[thick] (1.4,1.55)--(2.4,1.0);
 \draw[gray!65!black,dashed] (0.5,1)..controls(2.3,3.2)..(4.4,1.1);
 \node[font=\scriptsize,align=left] at (2.5,0.15){实线：局部邻域边\\虚线：不应直接使用的全局捷径};
\end{tikzpicture}
\caption{邻域图把局部观测关系变成边。图上的路径可以绕过折叠，直接连接会产生捷径。}
\label{fig:manifold-neighborhood-graph}
\end{figure}

\section{Isomap的全局测地距离}
\subsection{从邻域边到测地距离}
Isomap在2000年提出把非线性降维转化为图距离保持问题\cite{tenenbaum2000isomap}。其假设是流形上的短邻域边能够近似局部测地距离。对节点$i,j$，定义图最短路径距离
\begin{equation}
 d_G(i,j)=\min_{p:i\leadsto j}\sum_{(a,b)\in p}\ell_{ab}.
 \label{eq:isomap-graph-distance}
\end{equation}
$p$遍历图中从$i$到$j$的路径。连通时，该距离沿邻域图累加边长。折叠两段之间的直接欧氏距离不会成为候选边，除非邻域规则错误地把它们相连。

若图分成多个连通分量，跨分量的$d_G$没有有限定义，不能把无穷距离代入后续矩阵。增大邻域可能连接分量，却也可能引入捷径。实际计算应先检查连通性，而不是静默地用一个很大的常数替换无穷距离。

\subsection{经典多维尺度分析}
给定估计的距离$\delta_{ij}=d_G(i,j)$，Isomap使用经典多维尺度分析（Classical Multidimensional Scaling，cMDS）寻找$q$维坐标。令$D^{(2)}$满足$D^{(2)}_{ij}=\delta_{ij}^2$，中心化矩阵为
\begin{equation}
 H=I_n-\frac1n\bm1\bm1^{\mathsf T},
 \qquad B=-\frac12HD^{(2)}H .
 \label{eq:isomap-cmds}
\end{equation}
$\bm1$是$n$维全1向量，$H$去除坐标平移，$B$是由距离双中心化得到的内积矩阵。若距离来自欧氏点集，则$B$半正定。设$B=U\Lambda U^{\mathsf T}$，按正特征值降序排列，取前$q$个正特征值组成$U_q,\Lambda_q$，坐标为
\begin{equation}
 Z=U_q\Lambda_q^{1/2}.
 \label{eq:isomap-coordinates}
\end{equation}
因为$ZZ^{\mathsf T}=U_q\Lambda_qU_q^{\mathsf T}$，它是在低秩意义下逼近图距离所对应内积的坐标。若$B$出现显著负特征值，说明图距离未能形成可靠的欧氏几何或邻域图受到噪声和捷径影响。Isomap的全局距离保持目标对采样覆盖和图连通性敏感。

\section{LLE的局部重构坐标}
\subsection{局部权重的计算}
LLE同样在2000年提出，但它不要求恢复所有点对的全局距离\cite{roweis2000lle}。它只要求每个样本能由其近邻的局部仿射组合重构。令$\mathcal N_i$为样本$i$的$k$个近邻，权重$w_{ij}$只在$j\in\mathcal N_i$时允许非零，并满足
\begin{equation}
 \sum_{j\in\mathcal N_i}w_{ij}=1 .
 \label{eq:lle-affine-constraint}
\end{equation}
权重通过
\begin{equation}
 \min_{\{w_{ij}\}}\left\|\bm x_i-\sum_{j\in\mathcal N_i}w_{ij}\bm x_j\right\|_2^2
 \label{eq:lle-local-objective}
\end{equation}
求得。和为$1$使目标对平移不变。令$\bm r_{ij}=\bm x_i-\bm x_j$，则目标等于$\|\sum_jw_{ij}\bm r_{ij}\|_2^2$。若局部点近似位于切空间，权重主要由局部相对位置决定。

对邻域差向量组成矩阵$R_i=[\bm r_{ij}]_{j\in\mathcal N_i}\in\mathbb R^{d\times k}$，令$C_i=R_i^{\mathsf T}R_i$。若$\bm w_i$收集$k$个权重，则目标是$\bm w_i^{\mathsf T}C_i\bm w_i$。当$C_i$可逆时，拉格朗日方程给出
\begin{equation}
 \bm w_i=\frac{C_i^{-1}\bm1_k}{\bm1_k^{\mathsf T}C_i^{-1}\bm1_k}.
 \label{eq:lle-local-weights}
\end{equation}
其中$\bm1_k$为$k$维全1向量。推导来自$2C_i\bm w_i-\lambda\bm1_k=0$和$\bm1_k^{\mathsf T}\bm w_i=1$。当邻域点近似共面或$k>d$时，$C_i$常为奇异或病态。可用$C_i+\alpha\operatorname{tr}(C_i)I_k$替代$C_i$，其中$\alpha>0$是正则化系数，但这会改变局部重构关系。

\subsection{从重构保持到特征问题}
令$Y\in\mathbb R^{n\times q}$为低维坐标矩阵，并令$y_i^{\mathsf T}$为第$i$行。LLE要求原空间的重构权重在低维空间中仍然成立，因此最小化
\begin{equation}
 \Phi(Y)=\sum_i\left\|y_i-\sum_jw_{ij}y_j\right\|_2^2
 =\operatorname{tr}\left(Y^{\mathsf T}(I-W)^{\mathsf T}(I-W)Y\right).
 \label{eq:lle-embedding-objective}
\end{equation}
$W\in\mathbb R^{n\times n}$是稀疏权重矩阵。为排除所有坐标相同的零解，约束$Y^{\mathsf T}Y=I_q$并要求$Y^{\mathsf T}\bm1=0$。于是$Y$取矩阵$(I-W)^{\mathsf T}(I-W)$在常数向量之外最小的$q$个特征向量。局部权重若因噪声或邻域跨越边界而错误，低维坐标会继承该错误。LLE不保证全局距离、密度或外部标签被保持。

\section{图拉普拉斯与扩散距离}
\subsection{拉普拉斯特征映射}
拉普拉斯特征映射（Laplacian Eigenmaps）将“相邻点应有相近坐标”直接写成图平滑目标\cite{belkin2003laplacian}。给定对称权重$W=(w_{ij})$，定义度矩阵$D=\operatorname{diag}(d_1,\ldots,d_n)$，其中$d_i=\sum_jw_{ij}$。对一维坐标$f\in\mathbb R^n$，有
\begin{equation}
 \frac12\sum_{i,j}w_{ij}(f_i-f_j)^2=f^{\mathsf T}(D-W)f.
 \label{eq:laplacian-energy}
\end{equation}
左侧逐边惩罚相邻点的坐标差异。右侧定义图拉普拉斯矩阵$L=D-W$。取$\operatorname{diag}(f^{\mathsf T}Df)=1$并排除常数解，可求解广义特征问题
\begin{equation}
 Lf=\lambda Df.
 \label{eq:laplacian-eigenproblem}
\end{equation}
最小的非零特征向量提供低维坐标。若图断裂，$L$有多个零特征值，每个连通分量都贡献一个常数向量，坐标可能只是在编码分量，而非流形内部位置。

\subsection{扩散映射的多步邻域}
扩散映射把一次邻域相似度扩展为多步随机游走\cite{coifman2006diffusion}。令$P=D^{-1}W$，则$P_{ij}$是从节点$i$一步转移到$j$的概率。$P^t$给出$t$步转移概率。用合适的平稳分布$\pi$加权，扩散距离可写为
\begin{equation}
 D_t^2(i,j)=\sum_{k}\frac{(P^t_{ik}-P^t_{jk})^2}{\pi_k}.
 \label{eq:diffusion-distance}
\end{equation}
它比较两个节点经过$t$步后到达全图各位置的概率分布。短路边只影响局部转移，若其影响不足以改变多步分布，扩散距离会比单次欧氏距离稳定。$t$过小只看到噪声邻域，$t$过大则不同区域的分布可能都接近平稳分布，距离失去区分度。扩散坐标由$P$的非平凡特征向量按特征值的$t$次幂缩放得到。它因此同时受邻域尺度、权重核和时间尺度影响。

\section{概率邻域可视化}
\subsection{t-SNE的邻域概率目标}
t-SNE的目标是让高维空间中的近邻概率与低维空间中的近邻概率相匹配，而不是保持所有点对的距离\cite{vandermaaten2008tsne}。在高维空间，对$i\ne j$定义条件概率
\begin{equation}
 p_{j\mid i}=\frac{\exp(-\|\bm x_i-\bm x_j\|_2^2/2\sigma_i^2)}{\sum_{\ell\ne i}\exp(-\|\bm x_i-\bm x_\ell\|_2^2/2\sigma_i^2)}.
 \label{eq:tsne-high-prob}
\end{equation}
$\sigma_i$通常按给定困惑度选择，使每个点的有效邻域大小相近。对称概率为$p_{ij}=(p_{j\mid i}+p_{i\mid j})/(2n)$。低维点$y_i\in\mathbb R^q$使用重尾相似度
\begin{equation}
 q_{ij}=\frac{(1+\|y_i-y_j\|_2^2)^{-1}}{\sum_{a\ne b}(1+\|y_a-y_b\|_2^2)^{-1}},
 \qquad i\ne j .
 \label{eq:tsne-low-prob}
\end{equation}
优化两个分布的KL散度
\begin{equation}
 \operatorname{KL}(P\|Q)=\sum_{i\ne j}p_{ij}\log\frac{p_{ij}}{q_{ij}}.
 \label{eq:tsne-objective}
\end{equation}
该方向的KL惩罚会重点避免高维近邻在低维消失。它不对称地惩罚“近邻被拉远”，对远距离的精确保持较弱。不同随机初始化、困惑度和迭代设置可能得到不同图形。簇之间的相对距离、簇面积、坐标轴方向和簇大小不能直接解释为原空间的定量事实。t-SNE适合探索局部邻域，不适合未经检验的全局测量。

\subsection{UMAP的模糊邻域图}
UMAP以局部模糊邻域图为输入，再在低维空间优化相似度匹配\cite{mcinnes2018umap}。对每个点使用局部距离校正和尺度参数构造邻域隶属度，随后以模糊并集组合有向邻域。低维相似度通常取
\begin{equation}
 \widetilde w_{ij}=\left(1+a\|y_i-y_j\|_2^{2b}\right)^{-1},
 \qquad a,b>0,
 \label{eq:umap-low-similarity}
\end{equation}
并用交叉熵型目标奖励高维相邻点在低维相邻，惩罚低维中被错误连接的点对。参数$a,b$由目标的最小距离和曲线拟合确定，实际软件还涉及负采样和优化调度。UMAP通常比t-SNE更重视邻域图的连通结构，也可在降维模型上近似变换新样本。它仍然依赖距离、邻域数、最小距离和随机优化。图中的簇间距离与密度受这些选择影响，不能据此断言流形的真实体积或类别概率。

\section{方法选择与结果检查}

表\ref{tab:manifold-method-choice}按保持的对象比较本章方法。表中“全局距离”指点对距离的近似保持，“局部重构”指邻域权重关系，“图平滑”指相邻节点坐标差异，“可视化邻域”指显示高维近邻，不表示所有几何量都被保留。
\begin{table}[htbp]
\centering
\caption{按几何目标选择流形方法}
\label{tab:manifold-method-choice}
\renewcommand{\arraystretch}{1.25}
\begin{tabular}{>{\centering\arraybackslash}m{2.1cm}>{\centering\arraybackslash}m{2.3cm}>{\centering\arraybackslash}m{2.7cm}>{\centering\arraybackslash}m{2.1cm}}
\toprule
方法 & 主要保持对象 & 关键条件与风险 & 适合的第一问题 \\
\midrule
PCA & 线性投影与平方重构 & 弯曲结构会被拉直，尺度影响方差 & 数据是否接近线性子空间 \\
Isomap & 图上的全局测地距离 & 图需连通且少捷径，采样要覆盖流形 & 全局路径是否比欧氏距离合理 \\
LLE & 局部仿射重构权重 & 局部矩阵可能病态，邻域错误会传播 & 局部切空间是否稳定 \\
Laplacian Eigenmaps & 图上的局部平滑 & 图断连会产生多个零特征值 & 相邻状态是否应有相近坐标 \\
Diffusion Maps & 多步转移关系 & 时间尺度过小或过大会失去信息 & 哪些区域由多步随机游走相连 \\
\shortstack{t-SNE\\与UMAP} & 选定的邻域可视化 & 参数、随机性和密度解释有限 & 哪些局部群组值得进一步检查 \\
\bottomrule
\end{tabular}
\end{table}

评价流形表示应分成几何评价和任务评价。几何评价可检查邻域保持率、图连通分量、局部重构误差、成对距离相关性或不同随机种子的稳定性。任务评价应在训练、验证和测试数据上分别完成。若表示用于故障预测，应比较原始特征、PCA和流形表示在相同数据划分、相同信息可用性和相同分类器下的测试性能。标签可分性提高不代表因果机制被恢复。按转速或设备划分测试集，才能检查表示是否只编码了设备身份或运行条件。

\section{发展脉络与边界}

流形学习的早期里程碑形成了互补的机制链。Isomap在2000年把邻域图最短路径与经典多维尺度结合，说明局部边可以估计全局测地距离\cite{tenenbaum2000isomap}。LLE在同年用局部仿射重构权重替代全局距离\cite{roweis2000lle}。Laplacian Eigenmaps在2003年把图平滑写成广义特征问题\cite{belkin2003laplacian}。Diffusion Maps在2006年以多步随机游走定义扩散距离\cite{coifman2006diffusion}。t-SNE在2008年把邻域概率匹配用于二维和三维可视化\cite{vandermaaten2008tsne}。UMAP的预印本发表于2018年，将局部模糊邻域图与低维交叉熵优化结合\cite{mcinnes2018umap}。流形正则化在2006年把图平滑加入监督学习目标，使无标签样本也能约束函数在邻域图上的变化\cite{belkin2006regularization}。

Cunningham与Ghahramani的综述《Linear Dimensionality Reduction: Survey, Insights, and Generalizations》系统比较线性降维的投影、重构、概率和判别视角，可作为本章PCA基线的补充阅读\cite{cunningham2015linear}。Bi与Lafaye de Micheaux于2025年发表的预印本《A Survey and Comparative Evaluation of Intrinsic Dimension Estimators under the Manifold Hypothesis》专门比较内在维度估计器，并分析噪声、曲率、样本量和超参数选择对估计的影响\cite{bi2025intrinsic}。该综述关注流形维度估计，不覆盖本章介绍的全部降维算法，因此适合补充流形维度选择。

这些方法都需要把“邻近”当作可信信息。采样不覆盖区域、噪声制造近邻、图断连、尺度选择不当或数据实际包含离散跳变时，算法只能忠实地放大错误假设。流形维度估计也不是单一算法的输出标签。应结合局部谱、邻域稳定性、重构误差和任务测试判断。对于展示而非预测，图形应保留参数、随机种子、预处理和版本，避免把一次布局当成数据的唯一真相。

\section{小结}

本章从生成式流形假设出发，区分了观测维度、内在维度、局部坐标和噪声。邻域图把有限样本中的局部关系变成可计算对象。Isomap通过图最短路径近似测地距离。LLE保存局部重构权重。拉普拉斯特征映射保存图平滑，扩散映射比较多步随机游走。t-SNE和UMAP主要服务于邻域可视化，不能把图上的远近直接解释成原空间的物理距离。下一章将转向元学习，讨论如何在任务变化时利用已有经验，而不再把低维几何作为唯一问题。

\section*{阅读依据}
本章的经典来源包括Isomap、LLE、Laplacian Eigenmaps、Diffusion Maps、t-SNE、UMAP和流形正则化的原始论文\cite{tenenbaum2000isomap,roweis2000lle,belkin2003laplacian,coifman2006diffusion,vandermaaten2008tsne,mcinnes2018umap,belkin2006regularization}。Cunningham与Ghahramani的线性降维综述用于核对PCA基线的术语和目标。正文保留原始来源作为算法机制的直接依据。
